% !PRINTED 589
\weberpage{589}
%<<CONT
is. \gesperrt{We call $\Psi$ the group belonging to the ideal $\mathfrak{a}$, and we also say
that $\mathfrak{a}$ belongs to the group $\Psi$.}

Since in every equation between numbers and functionals in $\Omega$ one may carry out all
the substitutions of the group $\Phi$, whereby the quantities of the field $R$
remain unchanged, without the equation ceasing to hold, it follows that
units, integral functionals, associated functionals, and functionals divisible by one
another do not lose these properties when any one of the substitutions of
$\Phi$ is carried out. We single out the theorem:

1. \gesperrt{If $\mathfrak{p}$ is a prime ideal, then all the ideals $\mathfrak{p} \mid \varphi$
conjugate to $\mathfrak{p}$ with respect to $R$ are also prime ideals. If $\mathfrak{a}$
is divisible by any power of $\mathfrak{p}$ then $\mathfrak{a} \mid \varphi$ is divisible
by the same power of $\mathfrak{p} \mid \varphi$.} % [sic] German source omits the comma after "theilbar" and prints a stray comma after $\mathfrak p\mid\varphi$

If $\mathfrak{p}$ is a prime ideal in $\Omega$, there is one and only one natural prime
$p$ divisible by $\mathfrak{p}$.

If we decompose this into its prime factors in $R$, one of these prime factors, which
we denote by $\mathfrak{P}$, must be divisible by $\mathfrak{p}$.

If now $\mathfrak{A}$ is any ideal in $R$ which is at the same time divisible by $\mathfrak{p}$,
then the greatest common divisor of $\mathfrak{P}$ and
$\mathfrak{A}$, which by \S.\,139 is also contained in $R$, is divisible by $\mathfrak{p}$,
and is therefore certainly not a unit. $\mathfrak{P}$ and $\mathfrak{A}$ are thus
not relatively prime, and $\mathfrak{A}$ is consequently divisible by the prime ideal $\mathfrak{P}$.
If $\mathfrak{A}$ is also a prime ideal, then $\mathfrak{A}$ and
$\mathfrak{P}$ must be identical. With this, having regard to 1., it is proved:

2. \gesperrt{If $\mathfrak{p}$ is a prime ideal in $\Omega$, there is one and only one
prime ideal $\mathfrak{P}$ in $R$ that is divisible by $\mathfrak{p}$, and every ideal
in $R$ divisible by $\mathfrak{p}$ is at the same time divisible by $\mathfrak{P}$. The
prime ideal $\mathfrak{P}$ is at the same time divisible by all the
prime ideals $\mathfrak{p} \mid \varphi$ conjugate to $\mathfrak{p}$.}

The partial norm $\mathfrak{N}_R(\mathfrak{p})$ is divisible by $\mathfrak{p}$, and consequently, as
an ideal in $R$, by $\mathfrak{P}$. But by the theorem 1. it can, moreover, be
divisible by no ideal in $R$ different from $\mathfrak{P}$. For if $\mathfrak{P}'$
is any prime divisor of $\mathfrak{N}_R(\mathfrak{p})$, then $\mathfrak{P}'$ is divisible by
at least one of the prime factors $\mathfrak{p} \mid \varphi$ and is hence equal to $\mathfrak{P}$. Consequently
%CONT>>
